%!TEX root =  tfg.tex
\chapter{Metodología}

\begin{quotation}[Novelist]{Ernest Hemingway (1899--1961)}
The good parts of a book may be only something a writer is lucky enough to overhear or it may be the wreck of his whole damn life -- and one is as good as the other.
\end{quotation}

\begin{abstract}
En este capítulo se describe el marco de trabajo seleccionado para la ejecución del proyecto. Se detalla 
la estructura organizativa del equipo, el reparto de responsabilidades entre sus integrantes y la metodología 
ágil empleada. El objetivo es justificar cómo la combinación de técnicas de gestión (Kanban/Scrum) y 
herramientas de control de versiones garantiza un desarrollo ordenado, eficiente y capaz de adaptarse 
a los cambios técnicos surgidos durante la creación de Vía Inferni.
\end{abstract}

\section{Estructura organizacional del proyecto}

El proyecto ha sido desarrollado de forma conjunta por un equipo de dos integrantes, bajo un modelo de 
responsabilidad compartida y trabajo colaborativo. A diferencia de una estructura jerárquica rígida, 
se ha optado por un sistema de desarrollo transversal, donde ambos perfiles intervienen en todas las 
capas del software (arquitectura, gameplay y diseño) para garantizar la coherencia del producto y un 
aprendizaje integral de todas las áreas de la Ingeniería del Software.

Para organizar el trabajo de forma eficiente y evitar conflictos técnicos, las tareas se han distribuido 
siguiendo un esquema de especialización operativa, pero manteniendo una supervisión cruzada y constante:

\begin{itemize}
    \item \textbf{Co-desarrollo de Sistemas Críticos:} Ambos integrantes han participado activamente 
    en la definición de la arquitectura base (patrones de diseño y sistemas de eventos) y en el 
    núcleo del gameplay. Esto asegura que cualquier cambio en el motor de combate o en la generación 
    procedural sea conocido y validado por ambas partes, facilitando la integración de nuevos contenidos.
    
    \item \textbf{Gestión de Infraestructura Lógica y Datos (Integrante A):} Lidera la implementación 
    de la arquitectura desacoplada, el sistema de generación procedural de salas y la gestión de datos 
    mediante \textit{ScriptableObjects}. Para ello, cuenta con el apoyo constante del Integrante B en 
    la integración de los elementos visuales y el diseño de niveles dentro de estos sistemas lógicos.
    
    \item \textbf{Diseño de Comportamiento e Interfaz (Integrante B):} Coordina el desarrollo de la 
    inteligencia artificial de los enemigos y la lógica de la interfaz de usuario (HUD y menús). 
    En esta tarea, cuenta con el apoyo del Integrante A para asegurar que dichos subsistemas se 
    acoplen correctamente a la arquitectura global y a los sistemas de eventos del juego.
    
    \item \textbf{Validación y Pruebas:} Ambos integrantes comparten la responsabilidad de las 
    pruebas de integración y la corrección de errores críticos para asegurar que los módulos 
    individuales funcionen correctamente como un todo, garantizando la estabilidad de la demo técnica.
\end{itemize}

Ambos integrantes comparten la responsabilidad de las pruebas de integración y la corrección de errores 
críticos para asegurar que los módulos individuales funcionen correctamente como un todo, cumpliendo 
así con los estándares de calidad exigidos en un proyecto de Ingeniería del Software.

\section{Metodología de desarrollo}

Para el desarrollo de Vía Inferni se ha optado por una Metodología Ágil, específicamente un modelo 
híbrido entre Kanban y ceremonias ligeras de Scrum. Esta elección se debe a la necesidad de flexibilidad 
en un proyecto creativo donde las mecánicas de juego requieren iteración constante.

\subsection{Gestión de tareas con Kanban}

Dado que nuestro equipo es de dos personas, el uso de Git es crítico para evitar conflictos en los 
binarios de Unity:
\begin{itemize}
\item \textbf{Adaptación:} A diferencia de proyectos industriales con ciclos cerrados, nuestra adaptación 
de Kanban permite una gestión granular de tareas. Esto facilita que ambos integrantes puedan trabajar 
en módulos complementarios (por ejemplo, uno en la lógica de daño y otro en el HUD que la visualiza) 
sin bloquearse mutuamente, ajustando la prioridad de las tareas según las necesidades de la "build" actual.
\end{itemize}

\subsection{Sincronización y Comunicación (Scrum Ligero)}

Dado que nuestro equipo es de dos personas, el uso de Git es crítico para evitar conflictos en los 
binarios de Unity:
\begin{itemize}
\item \textbf{Sincronización diaria (Daily):} Reuniones breves de 15 minutos para poner en común el 
trabajo realizado, los objetivos de la jornada y, fundamentalmente, identificar bloqueos técnicos.
\item \textbf{Hitos de Sprint:} El calendario se ha dividido en fases críticas (Prototipado, Core Loop, 
Contenido y Pulido). Al final de cada fase, se realiza una evaluación del producto para decidir si una 
mecánica (como el cambio dual de Dante/Virgilio) funciona como se esperaba o requiere un rediseño antes 
de avanzar.
\end{itemize}

\subsection{Control de versiones y flujo de trabajo (GitFlow)}

Dado que nuestro equipo es de dos personas, el uso de Git es crítico para evitar conflictos en los 
binarios de Unity:
\begin{itemize}
\item \textbf{Ramas de funcionalidad (Feature Branches):} Se utiliza una rama independiente para cada 
nueva mecánica (ej: feature/generacion-procedural). Esto permite trabajar en paralelo en el sistema de 
niveles y en la IA sin riesgo de corromper la versión estable del juego.
\item \textbf{Revisión de Código:} Antes de fusionar cualquier funcionalidad en la rama de desarrollo 
(develop), el otro integrante debe revisar el código. Este proceso no solo actúa como filtro de calidad, 
sino que asegura que ambos comprendamos la lógica del compañero, manteniendo la transversalidad del equipo.
\end{itemize}

\subsection{Gestión de riesgos}

Se han identificado riesgos específicos del trabajo en equipo, como los conflictos en los prefabs de 
Unity o la posible desincronización de tareas. 
\begin{itemize}
\item \textbf{Mitigación:} Para mitigarlos, se ha establecido que las partes más complejas del código 
se traten mediante sesiones de Pair Programming (programación en pareja), asegurando que el conocimiento 
técnico sea compartido y no dependa de un solo integrante.
\end{itemize}
